%=====================================================================
%  main.tex — University of Gujrat SYNOPSIS (modular project)
%  Compile this file with pdfLaTeX (run TWICE for page references).
%
%  All formatting rules live in  uogformat.sty  — you should not need
%  to touch it. Edit only:
%    1. The details block just below (your name, reg #, supervisor...).
%    2. The files in  sections/  for your actual content.
%    3. references.tex  for your reference list.
%    4. The page numbers in toc.tex / listoffigures.tex /
%       listoftables.tex / listofappendices.tex after the final run.
%=====================================================================
\documentclass[12pt,a4paper]{article}
\usepackage{uogformat}

%---------------------------------------------------------------------
%  >>>  YOUR DETAILS  <<<  (used on the title page)
%---------------------------------------------------------------------
\newcommand{\SynopsisTitle}{UNIFIED MULTI-TASK DEEP LEARNING FRAMEWORK FOR PRECISION CATTLE MONITORING}
\newcommand{\ScholarName}{MUHAMMAD KHURAM}
\newcommand{\RollNo}{25014422-008}
\newcommand{\RegistrationNo}{25001222070}
\newcommand{\DegreeProgram}{MSc ELECTRICAL ENGINEERING}
\newcommand{\Department}{ELECTRICAL ENGINEERING}
\newcommand{\Faculty}{ENGINEERING}
\newcommand{\CampusName}{HAFIZ HAYAT CAMPUS}
\newcommand{\SupervisorName}{DR. SYED MUHAMMAD WASIF}
\newcommand{\EnrolmentSemester}{FALL 2025}
\newcommand{\FreezedSemesters}{NIL}
\newcommand{\SubmissionDate}{30-07-2026}
\newcommand{\DRRCApprovalDate}{15-08-2026}

\begin{document}

%---------------------- Preliminary pages: roman (i, ii, ...) --------
\pagenumbering{roman}
\pagestyle{prelims}

%=====================================================================
%  titlepage.tex — Synopsis title page (Rules 5.2, 6.1–6.5)
%  Edit your details in main.tex (the \newcommand block).
%=====================================================================
% Title page uses the SAME margins as the body (left 1.5in, right/top/bottom 1in),
% matching the actual UOG template file (its synopsis section, not the handbook's
% 1in front pages). No geometry override needed.
\thispagestyle{prelims}

% Title page uses TIGHTER 1.15 line spacing so all 11 fields + 5 signature
% lines fit on ONE page. (Body uses 1.5; change 1.15 below if you add/remove
% fields and it spills over.)
\begin{spacing}{1.15}

\begin{center}
% TITLE (Rule 6.2): CAPS, centered, bold, 16pt, with a rule line under it.
% \fontsize{16}{19}: 16 = size, 19 = line height. Edit 16 to resize the title.
{\fontsize{16}{19}\selectfont\bfseries \SynopsisTitle\par}
\rule{\textwidth}{1pt}                 % the horizontal bar under the title. 1pt = thickness.

\vspace{10pt}                          % gap below the bar (increase/decrease to taste)
% "BY" + SCHOLAR NAME (Rule 6.3): CAPS, centered, bold, 14pt.
{\fontsize{14}{17}\selectfont\bfseries BY\par}
\vspace{4pt}                           % gap between BY and the name
{\fontsize{14}{17}\selectfont\bfseries \ScholarName\par}
\end{center}

\vspace{4pt}                           % gap before the numbered details
% NUMBERED DETAILS 1-11 (Rule 6.4): bold, CAPS, left, 14pt.
% \fontsize{14}{18}: 14 = size, 18 = line height. parskip 0 = no extra gap between items.
{\fontsize{14}{18}\selectfont\bfseries\setlength{\parskip}{0pt}
% enumerate options:  label = "1." in bold;  leftmargin = indent of the list;
%   itemsep = vertical gap BETWEEN items (shrink to fit more);  topsep = gap above list.
\begin{enumerate}[label=\textbf{\arabic*.},leftmargin=2.2em,itemsep=3pt,topsep=0pt]
  \item ROLL No: \RollNo
  \item REGISTRATION No: \RegistrationNo
  \item DEGREE PROGRAM: \DegreeProgram
  \item DEPARTMENT: \Department
  \item FACULTY: \Faculty
  \item CAMPUS NAME: \CampusName
  \item SUPERVISOR NAME: \SupervisorName
  \item DEGREE ENROLMENT SEMESTER: \EnrolmentSemester
  \item FREEZED OR MISSES SEMESTER(S) (IF ANY): \FreezedSemesters
  \item DATE OF SYNOPSIS SUBMISSION TO THE DEPARTMENT: \SubmissionDate
  \item DATE OF APPROVAL FROM DRRC: \DRRCApprovalDate
\end{enumerate}
\par}

\vspace{12pt}                          % gap before signature block
% SIGNATURE LINES (Rule 6.5): bold, left, 14pt, NOT capitals.
% parskip 14pt = vertical gap BETWEEN signature lines (shrink to fit).
{\fontsize{14}{18}\selectfont\bfseries\setlength{\parskip}{14pt}
Scholar Signature: \dotfill           % \dotfill draws the "......" sign-here line

Supervisor Signature: \dotfill

Chairperson/Head Signature: \dotfill

Convener Faculty Board Signature: \dotfill

Date of approval by ASRB: \dotfill
\par}

\end{spacing}
\clearpage
          % (i)
%=====================================================================
%  toc.tex — Table of Contents (Rules 7.1, 7.2)
%  Tabular, horizontal borders only. Headings CAPS bold, entries normal.
%  >>> Update the page numbers (the \textbf{..} on the right) manually. <<<
%
%  Column spec  {@{}X@{\hspace{1.5em}}r@{\hspace{1em}}} means:
%    X            = Contents column (flexes to fill width; holds the dots)
%    \hspace{1.5em} = GAP between the dotted leader and the page number
%    r            = Page column, right-aligned
%    \hspace{1em} = right padding so numbers aren't jammed to the border
%  \specialrule{1.5pt}{0}{0} = thick (1.5pt) horizontal border.
%  \tocdots = the "......" leader (defined in uogformat.sty).
%=====================================================================
{\setlength{\parskip}{0pt}\begin{singlespace}
\noindent
\begin{tabularx}{\textwidth}{@{}X@{\hspace{1.5em}}r@{\hspace{1em}}}
\specialrule{1.5pt}{0pt}{0pt}
\multicolumn{2}{@{}c@{}}{\rule{0pt}{18pt}\makebox[\textwidth][c]{\textbf{Table of contents}}} \\[4pt]
\specialrule{1.5pt}{0pt}{0pt}
\multicolumn{2}{@{}c@{}}{\rule{0pt}{18pt}\makebox[\textwidth]{\hfill\textbf{Contents}\hfill\makebox[0pt][r]{\textbf{Page}}}} \\[4pt]
\hline
\rule{0pt}{16pt}%
LIST OF FIGURES\tocdots                          & \textbf{iii} \\[6pt]
LIST OF TABLES\tocdots                           & \textbf{iv}  \\[6pt]
LIST OF APPENDICES\tocdots                       & \textbf{v}   \\[6pt]
ABSTRACT\tocdots                                 & \textbf{01} \\[6pt]
INTRODUCTION\tocdots                             & \textbf{01} \\[6pt]
LITERATURE REVIEW\tocdots                        & \textbf{02} \\[6pt]
OBJECTIVES OF RESEARCH\tocdots                   & \textbf{02} \\[6pt]
RESEARCH QUESTION/PROBLEM\tocdots                & \textbf{03} \\[6pt]
METHODOLOGY (SAMPLING/MEASUREMENT/RESEARCH DESIGN/DATA ANALYSIS)\tocdots & \textbf{03} \\[6pt]
EXPECTED OUTCOMES OF RESEARCH\tocdots            & \textbf{04} \\[6pt]
REFERENCES\tocdots                               & \textbf{04} \\[6pt]
APPENDICES\tocdots                               & \textbf{05} \\[4pt]
\specialrule{1.5pt}{0pt}{0pt}
\end{tabularx}
\end{singlespace}\par}
\clearpage
                % (ii)
%=====================================================================
%  listoffigures.tex  (same tabular style as TOC — Rules 1.9.10 / 7.x)
%  Horizontal borders only, outer 1.5 pt. CAPS bold 11 pt heading.
%  Caption text in normal 11 pt with dashed leader to the page number.
%  >>> Add one row per figure and update page numbers manually. <<<
%=====================================================================
{\setlength{\parskip}{0pt}\begin{singlespace}
\noindent
\begin{tabularx}{\textwidth}{@{}X@{\hspace{1.5em}}r@{\hspace{1em}}}
\specialrule{1.5pt}{0pt}{0pt}
\multicolumn{2}{@{}c@{}}{\rule{0pt}{18pt}\makebox[\textwidth][c]{\textbf{LIST OF FIGURES}}} \\[4pt]
\specialrule{1.5pt}{0pt}{0pt}
\multicolumn{2}{@{}c@{}}{\rule{0pt}{18pt}\makebox[\textwidth]{\hfill\textbf{CONTENTS}\hfill\makebox[0pt][r]{\textbf{PAGE}}}} \\[4pt]
\hline
\rule{0pt}{16pt}%
Figure-6.1: Architecture of the unified multi-task framework
\leaders\hbox{-}\hfill\kern0pt & \textbf{03} \\[6pt]
\specialrule{1.5pt}{0pt}{0pt}
\end{tabularx}
\end{singlespace}\par}
\clearpage
      % (iii)
%=====================================================================
%  listoftables.tex  (same tabular style as TOC — Rules 1.9.11 / 7.x)
%  Horizontal borders only, outer 1.5 pt. CAPS bold 11 pt heading.
%  Caption text in normal 11 pt with dashed leader to the page number.
%  >>> Add one row per table and update page numbers manually. <<<
%=====================================================================
{\setlength{\parskip}{0pt}\begin{singlespace}
\noindent
\begin{tabularx}{\textwidth}{@{}X@{\hspace{1.5em}}r@{\hspace{1em}}}
\specialrule{1.5pt}{0pt}{0pt}
\multicolumn{2}{@{}c@{}}{\rule{0pt}{18pt}\makebox[\textwidth][c]{\textbf{LIST OF TABLES}}} \\[4pt]
\specialrule{1.5pt}{0pt}{0pt}
\multicolumn{2}{@{}c@{}}{\rule{0pt}{18pt}\makebox[\textwidth]{\hfill\textbf{CONTENTS}\hfill\makebox[0pt][r]{\textbf{PAGE}}}} \\[4pt]
\hline
\rule{0pt}{16pt}%
Table-6.1: Evaluation metrics and datasets planned\leaders\hbox{-}\hfill\kern0pt & \textbf{03} \\[6pt]
\specialrule{1.5pt}{0pt}{0pt}
\end{tabularx}
\end{singlespace}\par}
\clearpage
       % (iv)
%=====================================================================
%  listofappendices.tex  (Rules 1.9.12 / 7.x — same tabular style)
%=====================================================================
{\setlength{\parskip}{0pt}\begin{singlespace}
\noindent
\begin{tabularx}{\textwidth}{@{}X@{\hspace{1.5em}}r@{\hspace{1em}}}
\specialrule{1.5pt}{0pt}{0pt}
\multicolumn{2}{@{}c@{}}{\rule{0pt}{18pt}\makebox[\textwidth][c]{\textbf{LIST OF APPENDICES}}} \\[4pt]
\specialrule{1.5pt}{0pt}{0pt}
\multicolumn{2}{@{}c@{}}{\rule{0pt}{18pt}\makebox[\textwidth]{\hfill\textbf{CONTENTS}\hfill\makebox[0pt][r]{\textbf{PAGE}}}} \\[4pt]
\hline
\rule{0pt}{16pt}%
APPENDIX-01: Turnitin Originality Report
\leaders\hbox{-}\hfill\kern0pt & \textbf{05} \\[6pt]
\specialrule{1.5pt}{0pt}{0pt}
\end{tabularx}
\end{singlespace}\par}
\clearpage
   % (v)

%---------------------- Main body: arabic (Page X of Y) --------------
\clearpage
\pagenumbering{arabic}
\setcounter{page}{1}
\pagestyle{mainbody}

\section{Abstract}

Weed detection in precision agriculture is hindered due to need of tedious pixel-level annotation and domain shift between test and training data. A number of Visual Foundation Models (VFMs) such as the Segment Anything Model (SAM) have been used to solve the above problems. But its implementation often has a significant computational cost and is prone to semantic confusion between similar classes. In this thesis, a novel label-efficient framework for weed detection is proposed. It combines a lightweight few-shot object detector that will generate spatial prompts for VFMs. This will enable them to achieve high-fidelity zero-shot instance segmentation. Because it relies on an efficient zero-shot detector, it can be deployed quickly in the field without the need of exhaustive crop-specific retraining. The proposed architecture will be a computationally viable  universal perception system for scalable and autonomous precision weed management.
\section{INTRODUCTION}

The global yield loss caused by weeds amounts to 34\% on average (Oerke et al). The traditional homogeneous chemical weedicide application and labor-intensive hand weeding is becoming ineffective due to ecological impact, rising costs and the rapid evolution of herbicide-resistant weeds \cite{gao2025a}. There is a growing need for sustainable production methods. One such promising method is the development of automated precision agriculture solutions that use machine vision for site-specific precision weed management.

While  appealing in theory, deep learning-based approaches are currently limited by their reliance on fully supervised frameworks. These approaches require extensive pixel level annotated dataset, which is a time and resource-consuming process (up to 2.5 hours per image) \cite{nadeem2026}. Furthermore, the traditional computer vision models suffer from domain shift problem and fail catastrophically when deployed in a new field due to the significant variations in lighting, soil type and crop geometries \cite{magistri2023}. In this work, we propose a novel universal zero-shot weed detection framework hybrid foundation-model architecture, which requires only a small amount of annotated data.

\subsection{Background of the Study}

Recent advancements in agricultural computer vision reveal that one-stage detectors, like YOLO variants, effectively localize multiple targets at a fast speed in natural environments \cite{gao2024}. At the same time, Vision Foundation Models such as the Segment Anything Model (SAM) enable powerful zero-shot segmentation tasks through the use of a pre-trained architecture that does not require task-specific training data \cite{ghose2025}. Recent works highlight the possibility of combining such models in a cascade, where a fast, lightweight one-stage detector is used to provide spatial bounding box prompts to a foundation model, achieving superior performance over standard one-stage or two-stage detectors. However, such approaches are often memory-intensive and suffer from a computational complexity of quadratic complexity in the number of prompts, necessitating optimization in order to be deployed on agricultural hardware at the edge \cite{mo2025}.

\subsection{Scope of the Work}

This research will be limited to the development of a software which will be used to use the vision-based system to detect weeds on a robotic rover. The proposed approach will be based on the use of publicly available image sets of different crops (deng et al). The research effort will not include the development of the mechanical design of the rover, nor the development of physical mechanisms that will perform the weeding function.

\subsubsection{Assumptions}

The current study makes two critical assumptions in order to conduct the analysis. First, the data sets and the camera equipment used for this research will produce standard daylight optical images. Their quality and resolution are sufficient to allow for the generation of spatial prompts accurate enough for the foundational model \cite{ghose2025}. Second, the limitations inherent in the edge-computing hardware can be approximated in software via standard performance metrics, such as FPS and parameter count.
\section{LITERATURE REVIEW}

The development of robust, automated weed detection systems has been a central focus of precision agriculture research over the past decade. Making a universal weed detection framework is hindered by technical challenges, such as environmental variability, lack of annotated data, and the visual similarities between crops and weeds. This section critically examines the use of different computer vision techniques applied to this domain over the years. It progressed from early supervised machine learning approaches to the current state-of-the-art in artificial intelligence.

\subsection{Traditional Deep Learning in Precision Agriculture}
The shift from traditional machine vision to deep neural networks is one of the recent advances that increased the accuracy of automated weed recognition under controlled environments. Supervised Convolutional Neural Networks and one-stage object detectors, such as Yolo family networks, allow to automatically learn higher-level features from raw pixel data \cite{abedin2026} \cite{hamrani2025}. Recent developments in neural network design have shown outstanding performance in terms of speed and accuracy. According to \cite{sharma2024}, Yolo v11 achieved 13.5 ms inference speed and 89.10\% mAP on average. In turn, Yolov4-Tiny achieved 42.24 FPS on UAV maize weed detection \cite{pei2022}. The accuracy of a custom-made 5-layered Convolutional Neural Network reached 95.12\% with only 16.754 ms of inference speed and 0.012 GB memory consumption on a Raspberry Pi platform \cite{razfar2022}.

However, the accuracy of such methods is limited by the quality and quantity of annotated data, which is a known limitation of traditional computer vision (Deng et al., 2024). To train such a model, one needs to create pixel-level annotations, which is a very time-consuming process. According to (Gao et al., 2025), annotating one image could take 15-20 minutes. Therefore, it is not feasible to collect annotated data for every type of weed on the planet.

In addition to the lack of annotated data, deploying such a model in real-world conditions faces the problem of domain shift. In other words, if the distribution of data where the model is deployed differs from the distribution of data on which it was trained, its accuracy drops significantly (gao et al). The magnitude of this drop was recently quantified in an attempt to classify crops using a Source-Only approach: “On the cross-continental crop classification task, the mAP of a Source-Only model decreased by as much as 33.67\% from a value of 99.17\% (upper bound) to 65.50\% on the target set, completely failing to recognize the “Cotton” class” (Li et al., 2025). Similar results were presented when analyzing seasonal changes: “The baseline YOLOX, YOLOv8 and DINO detectors experienced performance drops ranging from 8.8\%, to 14.5\% on the cross-year experiment, with accuracy decreases of up to 26.1\% for DINO on the maize classification task” (Deng et al., 2024). Domain shift poses an even bigger challenge when it comes to recognizing plant species in the field. On images with a high degree of overlap and varying lighting conditions, “the mAP of Faster R-CNN fell to as low as 9.55\%” (Gao et al., 2025).


\subsection{Unsupervised Domain Adaptation (UDA) and Semi-Supervised Learning}
In order to save the astronomical amount of manual labor required to annotate a large number of images, and to reduce the catastrophic failures that occur due to domain shift, the agricultural research community has increasingly shifted to Unsupervised Domain Adaptation (UDA) and Semi-Supervised Learning (SSL) \cite{nadeem2026}. Initial methods based on Generative Adversarial Networks (GANs) and CycleGAN for style transfer produced results that were adequate for getting the source images translated to fit into target domains, with an Intersection over Union (IoU) increase of around 10\% on various source-target field pairs for the mean Intersection over Union (mIoU) metric \cite{ilyas2023} \cite{magistri2023}. To fill the significant visual gap between various plant images from the internet and the limited target domain agricultural robot data, Multi-level Attention-based Adversarial Discriminators (MAAD) have been incorporated into feature extractors, which resulted in 7.5\% higher object detection precision on unlabeled target domains \cite{tzouras2023}.

Likewise, semi-supervised teacher-student architectures have shown that using large amounts of unlabeled data can greatly improve accuracy. YOLOv8l-based WeedTeacher achieved a 2.6\% improvement over the fully supervised baseline with only 10\% of the labeled vegetable field (deng et al). However, other semi-supervised methods based on FCOS and Faster R-CNN on multi-class cotton datasets had about 76\% and 96\% detection accuracy, and greatly reduced the dependency on annotations (li et al). To avoid the model being overfitted by predicting false targets, class covariance analysis based pseudo labelling has been proposed \cite{huang2024}. New feature-level UDA improvements have enabled the development of masked adversarial alignment (CCUDA-SS) to selectively adapt models to uncertain areas while incurring minimal runtime complexity, resulting in up to +5.23\% mIoU improvement on eight datasets \cite{nadeem2026}. Moreover, 80.53\% mIoU was achieved by the 1-shot Supervised Domain Adaptation (SDA) model, which was very close to the 83.69\% accuracy of the fully supervised network when UDA was combined with the use of augmentation scheduling \cite{ilyas2023}.

But there are inherent limitations of UDA and SSL approaches, which limit their universal application. The most important weakness is "pseudolabel drift," which occurs when the domain changes significantly, and when a label is misclassified at the beginning, it is continually reinforced, actively contaminating the teacher model's target distribution \cite{nadeem2026}. Empirical results demonstrate that sub-optimal SSOD framework designs can negatively impact performance when domain shift is strong: for example, the DenseTeacher framework was applied to cross-domain weed datasets, yielding a significant decrease in mAP@50 of 13.4\% from DenseTeacher's fully supervised baseline (deng et al).


\subsection{The Paradigm Shift to Foundation Models and Zero-Shot Learning}
Looking beyond the iterative domain adaptation paradigm, the agricultural computer vision arena is witnessing a paradigm shift toward Visual Foundation Models (VFMs). After their introduction to the computer vision domain, vision transformers (ViTs) showed to have clear advantages over local CNNs in terms of capturing global relations \cite{dosovitskiy2020}. For instance, in a comparison between hierarchical ViTs (e.g. Swin Transformer) and standard ViT-B, the former achieved 92.3\% average accuracy on multi-crop settings \cite{ahmad2024}. As for novel architectures, the proposed WeedSwin model demonstrated outstanding performance for multi-stage classification tasks, attaining a test mAP of 0.993 at 218.27 FPS \cite{islam2025}. Meanwhile, to address the issue of large parameter budgets in most Transformers, a lightweight alternative was suggested, which slashed the number of parameters to 412K while managing to keep a reasonable inference time of 141ms for practical edge-level deployment \cite{yenna2026}.

Beyond the architecture design space, there has been an increasing number of works leveraging the powerful but simplistic Segment Anything Model (SAM) from Meta. For instance, the zero-shot adaptation ability of SAM was utilized for agricultural parcel boundary extraction in a phenology-aware manner \cite{simsek2026}. Furthermore, the same work used fine-tuned Stable Diffusion models to generate realistic synthetic imagery, thus minimizing the need for real image collection \cite{modak2024}. Due to their generic nature, foundation models lack the botanical knowledge to differentiate between crops and weeds; therefore, an interesting hybrid approach was introduced, which combined a lightweight object detector (YOLOv11) with SAM. Specifically, a few-shot fine-tuned YOLOv11 was utilized to generate bounding boxes, which, in turn, guided the segmentation performed by SAM2. This increased the mIOU by +11.3\% compared to a baseline while still maintaining a decent speed of 80.02 ms per image \cite{ghose2025}. In another hybrid SAM-YOLO approach, a Grounded YOLOv10-SAM2.1 model was used to segment cotton bolls, attaining an mIOU of 86.9\% at 18.9 ms per frame (approximately 52 FPS) (pooja et al).

Finally, while the VFMs demonstrated exceptional performance in many vision tasks, they show to have some shortcomings in the agriculture domain. For example, in a setting where the VFM is utilized as a zero-shot anomaly segmenter (i.e. consider weeds as anomalies), such models underperformed severely when facing extreme weed infestations. The reason for the failure was the skewed distribution of anomalies, which constituted the majority of the image in the most infested cases \cite{chong2025}. Lastly, a sanity check was done to see if foundational vision-language models such as Florence2 could be utilized for agriculture tasks without extra finetuning. However, the results showed that such an approach either underperformed significantly or performed worse than random (aashu et al).

\subsection{Multimodal Sensor Fusion and the Research Gap}
The repeated failures of both conventional CNNs and modern Foundation Models in highly occluded environments are dictated by the same limiting factor: the dependence on standard RGB cameras. Pure optical imagery is inherently limited in its ability to capture morphological details beyond the surface due to varying field illumination, shadows, and similar green hues of different plant species. An RGB-only segmentation model tested on the WeedsGalore dataset scored a dismal mIoU of 63.08\% (zeynep et al). Meanwhile, the fusion of different data streams poses its own challenges, including resolution mismatch, extreme temporal displacement between bands, and a lack of adaptive importance weighting (gao et al).

The way to overcome such limitations is to turn to alternative sources of information and attempt Multimodal sensor fusion. By combining the data from standard RGB, Near-Infrared (NIR), Red-Edge (RE), Thermal, or Short-Wave InfraRed (SWIR) sensors, one can obtain a more comprehensive understanding of the target scene. For instance, a simple fusion of RGB, NIR, and RE channels via a Transformer-CNN architecture allowed to boost the segmentation performance to 78.88\% mIUO with a mere 8.7 million parameters \cite{galymzhankyzy2025}. Another approach, implemented in the S3-Net framework, demonstrated the effectiveness of combining both RGB texture information and SWIR’s ability to peer through vegetation using Vision-Language Models. It achieved an incredible 96.9\% accuracy on seed phenotyping at 95 FPS (song et al). Despite the recent progress in self-supervised learning and contrastive vision models (MoCo v2, DenseCL), the majority of the studies show that self-supervised learning underperforms in comparison to conventional supervised training for plant phenotyping tasks. The reason for such a gap is tied to the ability of self-supervised models to capture spatial patterns, which is exacerbated by the high dimensionality and redundancy of plant imagery data \cite{ogidi2023}. To harness the power of such foundation models for plant phenotyping tasks, one has to implement aggressive model compression techniques. For instance, CDGLYOLOv11Lite utilizes operations fusion and TensorRT’s FP16/INT8 precision to achieve 67 FPS on a Jetson Xavier NX while only using 1.9 MB of memory \cite{mo2025}.

The analysis of the literature highlights the existence of a major Research Gap: while hybrid Foundation Models such as YOLOSAM address the issue of extensive manual annotation by enabling zero-shot segmentation, and Multimodal fusion helps to alleviate the ambiguity of optical data, the combination of the two has not been explored. In other words, there is little research on how to implement Multimodal fusion in a way that is computationally efficient enough for real-time edge processing. At the same time, most of the existing Foundation Models are optimized for unimodal RGB input. Thus, there is a need to develop a novel efficient multimodal foundation model that would be capable of zero-shot detection of various weed species in different crop fields.
\section{Objectives of Research}
The objectives of this research are as follows.

\begin{enumerate}[label=\arabic*.,leftmargin=2em,itemsep=2pt,topsep=0pt]

\item To develop a software framework that combines a few-shot object detector with a Visual Foundation Model to achieve high-fidelity, zero-shot pixel-level weed segmentation.

\item To systematically test the proposed modal framework across disparate unseen crop varieties and field environments without explicit retraining.

\item To evaluate and minimize the inference latency and memory footprint of on edge hardware.
\end{enumerate}

\section{Research Question / Problem}
Can a hybrid framework integrating a lightweight, few-shot object detector with a visual foundation model support zero-shot weed segmentation across diverse agricultural domains with accuracy comparable to fully supervised models, while reducing computational cost sufficiently for real-time edge deployment? Existing systems rely on domain-specific retraining and exhaustive manual annotations; the absence of a label-efficient, cross-domain foundation model architecture constitutes the research gap addressed here.
\section{Methodology (Sampling\slash Measurement\slash Research Design\slash Data Analysis)}
The proposed framework comprises a shared convolutional backbone with task-specific necks and three heads: a detection head, an ArcFace-based re-identification head, and a spatio-temporal behaviour head. Tracking is performed downstream using BoT-SORT-style association over the learned embeddings. The overall architecture is shown in Figure-\ref{fig:framework}.

% ---- FIGURE: caption BELOW, bold label "Figure-6.1:" (Rules 1.7.4–1.7.6) ----
% Helper macro: \uogfigure{file}{width}{caption}{label}
\uogfigure{framework.png}{0.85\textwidth}%
  {Architecture of the proposed unified multi-task framework (Reference if any)}%
  {fig:framework}

\subsection{Datasets}
Four public datasets will be used: MultiCamCows2024, Cows2021, CattleBehaviours6, and CVB. Since no single dataset provides all label types per frame, a masked-label training scheme activates only the losses for which ground truth exists in each batch.

\subsection{Evaluation Metrics}
Detection will be evaluated with mAP; re-identification with rank-1 accuracy and mAP; tracking with HOTA, IDF1, MOTA, and identity switches; behaviour recognition with top-1 accuracy and macro-F1. The plan is summarised in Table-\ref{tab:metrics}.

% ---- TABLE: caption ABOVE, bold label "Table-6.1:" (Rules 1.7.1–1.7.2, 1.7.6) ----
\begin{table}[ht]
  \caption{Evaluation metrics and datasets planned for each task of the proposed framework (Reference if any)}
  \label{tab:metrics}
  \begin{singlespace}
  \begin{tabularx}{\textwidth}{lXl}
    \specialrule{1.5pt}{0pt}{0pt}
    \textbf{Task} & \textbf{Metrics} & \textbf{Dataset} \\
    \hline
    Detection & mAP@0.5, mAP@0.5:0.95 & Cows2021 \\
    Re-identification & Rank-1, mAP & MultiCamCows2024 \\
    Tracking & HOTA, IDF1, MOTA, IDS & Cows2021 \\
    Behaviour recognition & Top-1 accuracy, macro-F1 & CattleBehaviours6, CVB \\
    \specialrule{1.5pt}{0pt}{0pt}
  \end{tabularx}
  \end{singlespace}
\end{table}

\subsection{Multi-Task Loss}
The framework is trained by minimising a weighted sum of the task losses, where the weights balance gradient magnitudes to mitigate negative transfer:
% ---- EQUATION: right justified, numbered in parentheses (Rule 1.8.1) ----
\begin{flushright}
$L_{total} = \lambda_{det} L_{det} + \lambda_{id} L_{arc} + \lambda_{beh} L_{beh}$ \hspace{2cm} (6.1)
\end{flushright}

\section{Expected Outcomes of Research}
The research is expected to deliver a jointly optimised, lightweight multi-task framework achieving accuracy comparable to task-specific pipelines at a fraction of the inference cost; empirical evidence on shared versus separate backbones (negative-transfer analysis); a masked-label training recipe for heterogeneous livestock datasets; and a deployable prototype validated on consumer-grade GPU hardware.


%=====================================================================
%  references.tex — References (Rules 9.1–9.3)
%
%  STYLE  : APA (Latest). Scholars of the Faculty of Engineering and
%           CS&IT may instead follow their faculty-approved style (IEEE).
%  FORMAT : font size 11, SINGLE line spacing, space before each entry
%           (the singlespace + parskip below handle this automatically).
%  ORDER  : APA — list alphabetically by first author's surname.
%           IEEE — list in citation order and number [1], [2], ...
%  RULE   : every reference here must be cited in the text.
%
%  Journal, book, and book-chapter entries use DIFFERENT patterns.
%  One model of each is given below — copy the matching one.
%=====================================================================
\begin{singlespace}
\setlength{\parskip}{12pt}


% The {99} tells LaTeX to expect up to 99 references so it aligns the numbers properly
\begin{thebibliography}{99}
	
	% --- Reference 1 ---
	% \bibitem{keyword} is what you use to link the text to this reference.
% 1. Abedin & Mehrub (First citation in 3.1)
\bibitem{abedin2026}
M. M.-H.-Z. Abedin and T. Mehrub, ``Comprehensive Review on Deep Learning and Foundation Models in Agriculture: Progress, Benchmarks, and Open Challenges,'' 2026.

% 2. YOLOv11 13.5ms inference (Anonymous)
\bibitem{sharma2024}
A. Sharma, V. Kumar, and L. Longchamps, ``Comparative performance of YOLOv8, YOLOv9, YOLOv10, YOLOv11 and Faster R-CNN models for detection of multiple weed species,'' 2024.

% 3. Raspberry Pi 5-layer CNN (Anonymous)
\bibitem{razfar2022}
N. Razfar et al., ``Weed detection in soybean crops using custom lightweight deep learning models,'' 2022.

% 4. Ghose et al. 
\bibitem{ghose2025}
P. Ghose et al., ``YOLO-SAM AgriScan: A Unified Framework for Ripe Strawberry Detection and Segmentation with Few-Shot and Zero-Shot Learning,'' 2025.

% 5. 15-20 minutes annotation time (Anonymous)
\bibitem{gao2025a}
X. Gao, J. Gao, and W. A. Qureshi, ``Applications, Trends, and Challenges of Precision Weed Control Technologies Based on Deep Learning and Machine Vision,'' 2025.

% 6. Li et al. (Cross-continental crop classification)
\bibitem{li2025}
S. Li et al., ``A Class-Aware Unsupervised Domain Adaptation Framework for Cross-Continental Crop Classification with Sentinel-2 Time Series,'' 2025.

% 7. Deng et al., 2024 (Cross-season detection / DINO drops)
\bibitem{deng2024}
B. Deng, Y. Lu, and J. Xu, ``Weed database development: An updated survey of public weed datasets and cross-season weed detection adaptation,'' 2024.

% 8. Faster R-CNN plummet to 9.55% (Anonymous)
\bibitem{gao2024}
J. Gao et al., ``Cross-domain transfer learning for weed segmentation and mapping in precision farming using ground and UAV images,'' 2024.

% 9. Nadeem et al., 2026 (CCUDA-SS)
\bibitem{nadeem2026}
N. Nadeem, M. H. Asad, and A. Bais, ``CCUDA-SS: Cross-Crop Unsupervised Domain Adaptation for Semantic Segmentation,'' 2026.

% 10. Ilyas et al., 2023 (Overcoming field variability)
\bibitem{ilyas2023}
T. Ilyas et al., ``Overcoming field variability: unsupervised domain adaptation for enhanced crop-weed recognition in diverse farmlands,'' 2023.

% 11. Deng et al., 2025 (WeedTeacher)
\bibitem{deng2025}
B. Deng, Y. Lu, and D. Brainard, ``Semi-Supervised Weed Detection in Vegetable Fields: In-domain and Cross-domain Experiments,'' 2025.

% 12. ViT parameters reduced 30% (Anonymous)
\bibitem{ahmad2024}
S. Ahmad et al., ``AI-Enabled Vision Transformer for Automated Weed Detection: Advancing Innovation in Agriculture,'' 2024.

% 13. Islam et al., 2025 (WeedSwin)
\bibitem{islam2025}
T. Islam et al., ``WeedSwin hierarchical vision transformer with SAM-2 for multi-stage weed detection and classification,'' 2025.

% 14. Grounded YOLOv10-SAM2.1 (Anonymous)
\bibitem{verma2026}
P. Verma and A. Paul, ``Toward Grounded YOLO-SAM: Unified Detection–Segmentation Framework for Agricultural Intelligence,'' 2026.

% 15. Chong et al., 2024 (Zero-Shot Semantic Segmentation)
\bibitem{chong2025}
Y. L. Chong et al., ``Zero-Shot Semantic Segmentation for Robots in Agriculture,'' 2025.

% 16. Soniya et al., 2024 (When YOLO meets SAM)
\bibitem{soniya2025}
Soniya, N. Ambadipudi, and A. Narayan, ``When YOLO Meets SAM: Data-Efficient Weed Density Estimation,'' 2025.

% 17. Reliance on standard RGB cameras (Anonymous)
\bibitem{hamrani2025}
A. Hamrani et al., ``AI and Robotics in Agriculture: A Systematic and Quantitative Review of Research Trends (2015–2025),'' 2025.

% 18. WeedsGalore dataset 63.08% (Anonymous)
\bibitem{galymzhankyzy2025}
Z. Galymzhankyzy and E. Martinson, ``Lightweight Multispectral Crop-Weed Segmentation for Precision Agriculture,'' 2025.

% 19. DeepLabv3+ massive 41.2 million params (Anonymous)
\bibitem{yenna2026}
G. R. Yenna et al., ``HybridViT-CAB: a vision transformer and convolutional attention network for precision weed detection in agricultural systems,'' 2026.

% 20. S3-Net framework (Song et al. / Shi et al.)
\bibitem{shi2026}
L. Shi et al., ``An AI-Driven Dual-Spectral Vision–Language Sensing Framework for Intelligent Agricultural Phenotyping,'' 2026.

% --- REMAINING UNUSED PAPERS (Sorted below so they don't mess up your sequence) ---

\bibitem{dosovitskiy2020}
A. Dosovitskiy et al., ``An Image is Worth 16x16 Words: Transformers for Image Recognition at Scale,'' 2020.

\bibitem{ogidi2023}
F. C. Ogidi, M. G. Eramian, and I. Stavness, ``Benchmarking Self-Supervised Contrastive Learning Methods for Image-Based Plant Phenotyping,'' 2023.

\bibitem{mo2025}
H. Mo et al., ``CDGLYOLOv11Lite: A Lightweight Field Maize–Seedling Weed Detector Based on Cross-Domain Shared Attention and a Conditional Gated Linear Unit,'' 2025.

\bibitem{luca2025a}
C. Luca, ``Comparison of Vision Transformers and CNNs in the Context of Agricultural Weed Detection,'' 2025.

\bibitem{tzouras2023}
V. Tzouras, L. Nalpantidis, and R. Güldenring, ``From Web Data to Real Fields: Low-Cost Unsupervised Domain Adaptation for Agricultural Robots,'' 2023.

\bibitem{magistri2023}
F. Magistri et al., ``From one field to another—Unsupervised domain adaptation for semantic segmentation in agricultural robotics,'' 2023.

\bibitem{modak2024}
S. Modak and A. Stein, ``Generative AI-based Pipeline Architecture for Increasing Training Efficiency in Intelligent Weed Control Systems,'' 2024.

\bibitem{li2024}
J. Li et al., ``Performance evaluation of semi-supervised learning frameworks for multi-class weed detection,'' 2024.

\bibitem{simsek2026}
F. F. Şimşek and M. Altay, ``Phenology aware agricultural boundary extraction using segment anything model and planet scope imagery (zero shot learning approach),'' 2026.

\bibitem{luca2025b}
C. Luca, ``Transfer Learning and Fine-Tuning Vision Transformers for Multi-Crop Weed Classification,'' 2025.

\bibitem{huang2024}
Y. Huang and A. Bais, ``Unsupervised Domain Adaptation for Weed Segmentation Using Greedy Pseudo-labelling,'' 2024.

\bibitem{pei2022}
H. Pei et al., ``Weed Detection in Maize Fields by UAV Images Based on Crop Row Preprocessing and Improved YOLOv4,'' 2022.

	
% ---------- APA: Journal article ----------
% Author, A. A., & Author, B. B. (Year). Title of article. Journal
% Name, Volume(Issue), pages.


% ---------- APA: Book ----------
% Author, A. A. (Year). Title of the book (Edition). Publisher.


% ---------- APA: Book chapter ----------
% Author, A. A. (Year). Title of chapter. In E. Editor (Ed.), Title of
% book (pp. xx–xx). Publisher.


% ---------- IEEE: Conference paper (Engineering / CS&IT) ----------
% [n] A. Author, B. Author, and C. Author, ``Paper title,'' in Proc.
%     Conf. Name (ABBR), Year, pp. xx–xx.
% [1] Z. Wang, L. Zheng, Y. Liu, Y. Li, and S. Wang, ``Towards real-time
%     multi-object tracking,'' in \textit{Proc. ECCV}, 2020, pp. 107–122.
\end{thebibliography}
\end{singlespace}

%=====================================================================
%  appendices.tex — Appendices (Rules 10.1, 10.2)
%=====================================================================
\section{Appendices}

\clearpage
\begin{flushright}
{\fontsize{16}{19}\selectfont\bfseries\underline{APPENDIX-01}}
\end{flushright}

Print of system generated Turnitin similarity report (only `originality report' showing `primary sources') \textbf{signed by the author and the supervisor} is required to be placed here.

\begin{center}
The last page of such report shows `excluded items'.
\end{center}


\end{document}
